\section{Risks, kill criteria and pivots}
\label{sec:risks}

Likelihood and impact are qualitative (L/M/H), because nothing in the validation program
(\cref{sec:validation}) has run yet to give a measured rate. Each kill or pivot criterion is either
already registered (\emph{REORIENTATION.md} \S22) or proposed here for pre-registration before the
relevant gold or human data are seen. The external gate reader applies the criteria at MS1, MS3
and MS5 (\cref{sec:gate-reader}). \Cref{tab:risks} lists the risks that can stop or redirect the
program. \evfut{} Every row describes a future event, not an observed one.

\begingroup
\footnotesize
\begin{longtable}{@{}p{0.14\linewidth}p{0.05\linewidth}p{0.24\linewidth}p{0.19\linewidth}p{0.21\linewidth}@{}}
  \caption{Program risks, kill criteria and pivots. \enquote{Kill} stops a claim and the money
    that depends on it; \enquote{pivot} redirects the program to a narrower claim on its own,
    separate budget.}
  \label{tab:risks} \\
  \toprule
  Risk & L/I & Early signal & Mitigation & Kill / pivot criterion \\
  \midrule
  \endfirsthead
  \multicolumn{5}{@{}l}{\small\itshape \tablename~\thetable{} (continued)} \\
  \toprule
  Risk & L/I & Early signal & Mitigation & Kill / pivot criterion \\
  \midrule
  \endhead
  \midrule
  \multicolumn{5}{@{}r}{\small\itshape continued on next page} \\
  \endfoot
  \bottomrule
  \endlastfoot
    \textbf{Closure-step failure}: the absence check cannot be made informative &
      H/H &
      \evobs{} On every existing ledger the closure judge's rate on a candidate's own evidence
      equals its rate on swapped evidence (\cref{sec:n3}) &
      Phase~0 runs first: a written codebook with anchors per lens, a non-owner blind coder, a
      corpus-level recall check, and a judge panel including other-family and frontier judges &
      \textbf{MS1}: human $\kappa<0.70$, recall below its floor, or no judge beating the
      mismatched-evidence control by the margin \(\to\) stop the automated absence check; no grant is requested for
      the gap map \\
    \textbf{Central scientific kill}: the gap map does not beat the strongest baseline &
      M/H &
      A near-zero pooled point estimate; the PLC ranking already tracks source density
      (Spearman~$\rho=0.91$, \evobs{})\footnote{\repo{research/n3/plc/gapmap.json}.} &
      Density and area size fixed as baselines in advance; an incremental analysis over both;
      the strongest baseline chosen inside each bootstrap resample; $\delta$ fixed before any
      gold &
      \textbf{MS3} fails its screening rule \(\to\) no Phase~2 money; the interview-only pivot on
      its own budget. \textbf{MS5}: anything short of continue \(\to\) H2 reported as not
      supported, and H2 funding ends \\
    \textbf{Gold unavailable or too few golds} &
      M/H &
      The Phase~1 metadata check finds fewer than three eligible itemized golds, or golds that
      cover under 15 missed areas &
      Candidate golds listed before Phase~1 starts; pooling by random effects across all golds
      that pass &
      Fewer golds are reported at MS2; the MS3 rule does not change, so CONTINUE becomes harder
      to reach and Phase~2 is likelier to stay unreleased. No other study stands in for the
      missing golds \\
    \textbf{Contamination / memorization} &
      H/M &
      A high share of probe-positive items; recall on an old, widely cited gold far above what
      post-cutoff items show &
      Dated corpora from archived or publication-dated sources; per-item memorization probes with
      exclusion \parencite{golchin2023time,sainz2023nlp}; open-weight models with documented
      cutoffs &
      Probe-positive items dominate a gold \(\to\) drop that gold from the pooled reading and
      report its recall only as an upper bound \\
    \textbf{Expert recruitment} (V5) &
      M/H &
      Slow enrollment or no-shows in the V5 pilot; the pilot projects more areas than the funded ceiling
      of 450 &
      Professional bodies and honoraria; sessions capped at 75 minutes; about 20 short probes per
      expert; three experts per area &
      \textbf{MS4}: the simulated size exceeds the ceiling \(\to\) V5 main is not released, and
      H2 is reported as untested prospectively \\
    \textbf{Governance}: the owner overrides a gate &
      M/H &
      An exception taken by the owner, as in the step from N2 to N3, which no independent person
      checked (\cref{sec:evolution}) &
      The external gate reader reads MS1, MS3 and MS5 and must sign off any owner exception
      (\cref{sec:gate-reader}) &
      An exception without the reader's sign-off voids that gate reading; no tranche is released
      on it \\
    \textbf{Team / key-person risk} &
      H/H &
      No host and no named team yet; one person built, ran, coded and read the PoC
      (\cref{sec:team}) &
      Named host, PI, statistician and CTA collaborator before submission; non-owner coders for
      every gating score; a second engineer in Phase~1 &
      The key person is lost \(\to\) the program pauses at its last gate; no gate reading is
      accepted without the gate reader \\
    \textbf{H-OSS fails} (organizational track) &
      M/M &
      Too few usable departures in V3's pilot; the code lenses fire mostly on commit style &
      Code lenses frozen at MS2b, after the Git adapter exists; knowledge-at-risk and authorship
      baselines; a power decision after the pilot &
      The H-OSS gate reads stop \(\to\) H-OSS is reported as not supported. H2 is unaffected
      either way \\
    \textbf{Partner access} (V7) &
      M/H &
      No letter of intent by the pilot's scouting deadline &
      Scouting starts only after MS5; a software organization first; an artifact-coverage-only
      offer that never evaluates named individuals &
      No partner \(\to\) report the organizational-transfer claim as untested \\
    \textbf{Ethics / legal} (data protection, works councils) &
      M/H &
      A data protection impact assessment finds high residual risk, or a works council objects &
      Artifact-coverage mode by default; pseudonymized concentration scores; no evaluation use of
      outputs, by contract; on-prem processing where required &
      The assessment or the works council blocks the design \(\to\) stop the pilot at that
      partner; reconsider the product design before S6 \\
    \textbf{Cost / budget headroom} &
      L--M/M &
      A provider spending limit sits below the sum of a stage's own registered caps, as it did in
      the PoC\footnote{\evobs{} \repo{docs/lessons.md}: a \$5 weekly key limit stopped a run at
      its first call; the registered caps summed to about \$21.} &
      Provider limits at least 3$\times$ the registered caps, on a monthly window; the headroom
      check runs as code; one key per experiment (\cref{sec:resources}) &
      A study stopped by budget is reported as inconclusive and rerun with corrected headroom; a
      budget stop is never read as a negative result \\
\end{longtable}
\endgroup

Three further risks recur across the studies. They have no kill rule of their own, but they shape
how any result above should be read.

\textbf{In-sample tuning.} The gap map's lexicons were tuned on the same three ledgers the PoC
used to build it, so any effect measured on those ledgers is optimistic. The configuration is
hash-frozen before any gold or new domain is touched. Any change after the freeze turns a result
into an exploratory one, not a gate reading.

\textbf{Retrieval instability.} Two independently reconstructed ledgers on the same regulation
shared only 3 of 13 and 3 of 9 unclosed lens records (gap candidates plus single-source retrieval
gaps, \cref{sec:n3}). Low cross-run agreement under a frozen corpus would mean the gap map is not
reproducible. This is checked before any human time is spent on it.

\textbf{Surveillance framing.} A partner's staff can read a knowledge-coverage score as covert
performance monitoring, whether or not the score is accurate. The mitigation is contractual:
aggregate output by default, no evaluation use, and consent for naming. A partner that uses the
output to evaluate a named person is grounds to end the pilot (\cref{sec:organizations}).

\subsection{Pivots, and what each one keeps}

Two pivots exist. Each has its own budget line, and neither inherits Phase~2's money
(\cref{sec:resources}).

\textbf{The interview-only pivot.} If MS1 stops the automated absence check, or MS3 or MS5 ends
H2 funding, the program can propose one smaller study. It compares question selection by expected
information gain with untargeted interviewing, with no reconstruction prior. It keeps a light
workbench, the blind-arm protocol, the coding tools and the equivalence-test machinery; only the
reconstruction pipeline's role as predictor is dropped. A positive result is a usable efficiency
tool. A negative one bounds how far text reconstruction can go, a bound the field currently lacks
(\cref{sec:contributions}).

\textbf{Reconstruction with provenance as a documentation-audit tool.} If the human-facing claim
fails but the ledger and its provenance checks still work, the pipeline keeps a use that does not
depend on predicting hidden knowledge. It becomes a provenance-checked map of what a text states
and does not state, reported as a documentation audit, not a residual predictor. \evint{} This also
bears on a commercial risk: a vendor may fold a similar coverage score into its platform first.
The program's answer is not a race to ship. It is the released closure set, the benchmark, and a
prospective $\Delta$AUROC if MS5 ever reads continue (\cref{sec:contributions}).
